\section{Randomized threshold batches close the arbitrary-graph word-model gap}
\label{sec:two-stage-randomized-op2}

The preceding sections were written before the companion
\texttt{active\_edge\_lcp} note proved OP2.  Its decisive idea is not a new
terminal linear solver.  It is a square-root bound on the number of
\emph{complete exposed faces} that must be solved before the unreleased RPPR
energy is negligible.  This section records the exact transfer and corrects
the earlier statement that arbitrary-graph accelerated discovery was wholly
open.

Put
\begin{equation*}
 c_\rho=b-\alpha\rho D^{1/2}\one,
 \qquad
 F_\rho(x)=\frac12x^\top Qx-c_\rho^\top x
 \quad(x\geq0).
\end{equation*}
For a reachable active face $U$, let
\begin{equation}
 x_i^U=
 \begin{cases}
  (Q_{UU}^{-1}c_{\rho,U})_i,&i\in U,\\
  0,&i\notin U,
 \end{cases},
 \label{eq:two-stage-random-face}
\end{equation}
where reachability means that the seed is admitted first and every later
batch has strictly positive exact residual $c_\rho-Qx^U$.  Stieltjes inverse
positivity gives
\begin{equation}
 U\subseteq S^\star,
 \qquad
 0<x_U^U\leq x_{\rho,U}^\star.
 \label{eq:two-stage-random-safe-face}
\end{equation}
The trivial regime $\rho\geq1/d_v$ has $x_\rho^\star=0$ and is recognized by
one seed-degree query; all nonempty faces below refer to $\rho<1/d_v$.

We import the following proved theorem from the companion note.  It is stated
here in the exact form needed by the two-stage composition.

\begin{theorem}[Threshold-batch source theorem]
\label{thm:two-stage-threshold-source}
In the canonical nonzero regime, fix an RPPR objective target
$0<\varepsilon_{\rm obj}<\alpha/2$ and a
failure probability $\zeta\in(0,1)$.  There is an exact-real randomized local
algorithm which returns an explicit failure flag with probability at most
$\zeta$ and otherwise returns a reachable face $U$ and a residual-certified
face point $z$.  Its exact face satisfies
\begin{equation}
 F_\rho(x^U)-F_\rho(x_\rho^\star)\leq\varepsilon_{\rm obj},
 \label{eq:two-stage-threshold-face-gap}
\end{equation}
and the projected numerical point satisfies the same target.  Every active
face is contained in $S^\star$, hence
$\vol(U)\leq1/\rho$, and the fully charged expected work is
\begin{equation}
 \widetilde O\!\left(
   \frac{1}{\rho\sqrt\alpha}
   \log\frac1{\varepsilon_{\rm obj}}
   \log\frac2\zeta
 \right).
 \label{eq:two-stage-threshold-source-work}
\end{equation}
No support, ambient preprocessing, boundary oracle, or complementarity
margin is supplied.
For $\varepsilon_{\rm obj}\geq\alpha/2$, zero already meets the objective
target; the seed-degree test handles the trivial zero-RPPR regime.
\end{theorem}

The construction sets
\begin{equation}
 \vartheta=\frac18\sqrt{\alpha\rho\varepsilon_{\rm obj}},
 \qquad \nu=\vartheta/4,
 \qquad
 J=\left\lceil
  \frac{\log(64/\varepsilon_{\rm obj})}
       {2\log(1/q_\alpha)}
 \right\rceil,
 \qquad
 q_\alpha=\frac{\sqrt{2/\alpha}-1}{\sqrt{2/\alpha}+1}.
 \label{eq:two-stage-threshold-parameters}
\end{equation}
At each phase it constructs the exposed degree-coordinate SDD matrix, makes
capped independent calls to the randomized solver of
\citet{koutis2011nearly}, and accepts a call only if a complete active-row
scan certifies
\begin{equation}
 \norm{Q_{UU}z_U-c_{\rho,U}}_2\leq\nu\sqrt\alpha.
 \label{eq:two-stage-threshold-residual-test}
\end{equation}
It then admits every boundary row whose approximate positive residual is
larger than $\vartheta/2$.  The residual test makes every admission
deterministically support-safe and leaves each unreported exact residual
below $\vartheta$.

For completeness, the source of the square-root phase count is a block
Cholesky argument.  Order the final support by admission batches, retain only
the diagonal and first block subdiagonal of its scalar Cholesky factor, and
call the resulting block-bidiagonal M-matrix $\widehat C$.  Entrywise inverse
ordering and the two-block row structure give
\begin{equation}
 \norm{\widehat C^{-1}}_2\leq\frac1{\sqrt\alpha},
 \qquad
 \norm{\widehat C}_2\leq\sqrt2.
 \label{eq:two-stage-bidiagonal-conditioning}
\end{equation}
Chebyshev inverse decay along this block chain proves
\begin{equation}
 F_\rho(x^{U_J})-F_\rho(x_\rho^\star)
 \leq 8q_\alpha^{2J}+\frac{\vartheta^2}{\alpha\rho}.
 \label{eq:two-stage-batch-depth-import}
\end{equation}
The distributed threshold term charges every delayed row once, at its
eventual elimination block.  Since
$\log(1/q_\alpha)\geq\sqrt{2\alpha}$, only
$O(\alpha^{-1/2}\log(1/\varepsilon_{\rm obj}))$ complete faces are rebuilt.
The full factor proof and randomized-solver ledger remain in the companion
audit rather than being repeated here.

The new observation is that the face in
\cref{thm:two-stage-threshold-source} need not contain $S^\star$.

\begin{lemma}[Face energy is an approximate-envelope certificate]
\label{lem:two-stage-inner-face-energy}
Let $U$ be reachable and let
\begin{equation}
 \Gamma_U:=F_\rho(x^U)-F_\rho(x_\rho^\star).
 \label{eq:two-stage-inner-face-gap}
\end{equation}
Then the exterior RPPR amplitude and the ordinary principal-PPR solution
$u_U=Q_{UU}^{-1}b_U$ obey
\begin{align}
 \max_{i\notin U}\frac{x_{\rho,i}^\star}{\sqrt{d_i}}
 &\leq\sqrt{\frac{2\Gamma_U}{\alpha}},
 \label{eq:two-stage-inner-amplitude}\\
 \norm{D^{-1/2}(x^0-\widetilde u_U)}_\infty
 &\leq\rho+\sqrt{\frac{2\Gamma_U}{\alpha}}.
 \label{eq:two-stage-inner-ppr-error}
\end{align}
Thus a certified low-energy \emph{inner} face is a valid Stage~I set
certificate even when $U\subsetneq S^\star$.
\end{lemma}

\begin{proof}
Strong convexity of the RPPR objective on the nonnegative orthant gives
\begin{equation*}
 \Gamma_U\geq\frac\alpha2\norm{x^U-x_\rho^\star}_2^2.
\end{equation*}
The zero extension of $x^U$ vanishes outside $U$, and every $d_i\geq1$.
Hence each omitted normalized coordinate is at most
$\sqrt{2\Gamma_U/\alpha}$, proving
\cref{eq:two-stage-inner-amplitude}.  Apply the approximate-envelope
linearization proposition with
$\delta_U=\max_{i\notin U}x_{\rho,i}^\star/\sqrt{d_i}$ to obtain
\cref{eq:two-stage-inner-ppr-error}.
\end{proof}

\begin{theorem}[Randomized literal two-stage completion]
\label{thm:two-stage-randomized-completion}
For the nontrivial regime $0<\epsppr<1$, choose positive budgets
$\rho,\eta,\tau$ satisfying
\begin{equation}
 \rho+\eta+\tau\leq\epsppr,
 \qquad
 \varepsilon_{\rm obj}=\frac{\alpha\eta^2}{2}.
 \label{eq:two-stage-random-budget}
\end{equation}
If the seed-degree test certifies the zero-RPPR regime, return zero and omit
Stage~II; the RPPR--PPR bridge gives error at most $\rho$.
Run \cref{thm:two-stage-threshold-source} with failure budget $\zeta/2$ as
Stage~I, retain only its face $U$, and discard its numerical point and
randomized solver state.  In
Stage~II solve the ordinary principal system
\begin{equation}
 Q_{UU}u_U=b_U
 \label{eq:two-stage-random-tail}
\end{equation}
to Euclidean error at most $\tau$ with failure budget $\zeta/2$, take its
positive part, and then zero-pad the result.  With failure probability at
most $\zeta$, the output is semantically
$\epsppr$-accurate.  Using a second residual-certified randomized SDD solve
for Stage~II, the expected fully charged work is
\begin{equation}
 \widetilde O\!\left(
  \frac{1}{\rho\sqrt\alpha}
  \log\frac1{\alpha\eta^2}
  \log\frac4\zeta
  +\frac1\rho
  \log\frac1{\sqrt\alpha\tau}
  \log\frac4\zeta
 \right).
 \label{eq:two-stage-random-total-work}
\end{equation}
In particular, $\rho=\eta=\tau=\epsppr/3$ gives
$\widetilde O(1/(\epsppr\sqrt\alpha))$ work on every graph in the stated
exact-real randomized model.
\end{theorem}

\begin{proof}
On a successful Stage~I execution,
\cref{thm:two-stage-threshold-source,lem:two-stage-inner-face-energy} give
\begin{equation*}
 \norm{D^{-1/2}(x^0-\widetilde u_U)}_\infty
 \leq\rho+\eta.
\end{equation*}
The terminal Euclidean error contributes at most $\tau$ to the normalized
$\ell_\infty$ error because $d_i\geq1$.  Euclidean projection onto the
nonnegative orthant cannot increase distance from the exact nonnegative
principal solution.

For the Stage~II work claim, use the degree-coordinate SDD matrix
$D_U^{1/2}Q_{UU}D_U^{1/2}$.  The exact principal PPR solution has energy at
most $\alpha$, so a good randomized call at relative energy tolerance
$\sqrt\alpha\tau$ has residual norm at most $\alpha\tau$.  Accept a call
only after an exact active-row residual scan verifies that inequality.  Every
accepted call then has Euclidean solution error at most $\tau$,
deterministically.  Capped independent retries give the declared failure
probability.  The face has volume at most $1/\rho$; applying the source SDD
work theorem and adding \cref{eq:two-stage-threshold-source-work} proves
\cref{eq:two-stage-random-total-work}.
\end{proof}

\begin{corollary}[The shortest algorithm returns after Stage I]
\label{cor:two-stage-randomized-direct-return}
Set
\begin{equation*}
 \rho=\epsppr/2,
 \qquad
 \varepsilon_{\rm obj}=\alpha\epsppr^2/8.
\end{equation*}
The projected numerical point already returned by
\cref{thm:two-stage-threshold-source} is semantically
$\epsppr$-accurate, with expected work
$\widetilde O(1/(\epsppr\sqrt\alpha))$.  Therefore the literal Stage~II in
\cref{thm:two-stage-randomized-completion} is a modular set-only variant, not
the fastest required execution.
\end{corollary}

\begin{proof}
The standard RPPR objective conversion gives normalized PPR error at most
\begin{equation*}
 \rho+\sqrt{2\varepsilon_{\rm obj}/\alpha}=\epsppr.
\end{equation*}
Substitute $\rho=\epsppr/2$ in
\cref{eq:two-stage-threshold-source-work}.
\end{proof}

\begin{remark}[Las Vegas wrapper]
All accepted face solves and admissions are verified deterministically; the
randomness can only cause an explicit failure flag.  Taking a constant
$\zeta<1$ and restarting the complete procedure after a flag therefore gives
a Las Vegas exact-real algorithm with constant expected restart overhead.
This does not supply deterministic bit complexity.
\end{remark}

\paragraph{What transfers to AESP--CD.}
The safe-residual admission rule, hard face cap, restart discipline, and
certificate race transfer directly.  The result should be viewed as a new
threshold-batched active-set lane compatible with the AESP--CD portfolio, not
as a proof that the original momentum recurrence itself has become
randomized.  It deliberately rebuilds a fresh exposed SDD solver on each of
only $\widetilde O(1/\sqrt\alpha)$ faces; no momentum, Perron alignment,
face-shock bank, or dynamic response structure crosses a batch.

\paragraph{Remaining scope.}
The arbitrary-graph existence question is closed only in the companion
note's exact-real randomized word model.  Deterministic finite-precision and
coefficient-bit bounds, a practical local implementation of the randomized
SDD primitive, and a faster persistent changing-face response remain open.
The earlier Green, leakage, APPR, SOR, tree-message, and mass lanes remain
valuable as deterministic or practical competitors; they are no longer the
only route to the target asymptotic existence theorem.
